% ==============================================================================
% SLIDE 4: URBAN MORPHOLOGY & DEFENSIVE VERNACULAR
% ==============================================================================
\begin{frame}[label=defense]{Urban Morphology: Wind, Water \& Defense}
  \begin{columns}[T, onlytextwidth]
    % Left Column: Photo / Streetscape
    \begin{column}{0.48\textwidth}
      \slidephoto{Mykonos-street.JPG}{Dense urban fabric and labyrinthine pathways in Mykonos}
    \end{column}

    % Right Column: Bullet Points
    \begin{column}{0.48\textwidth}
      \textbf{\color{aegeanblue}Urban Fabric as Defense System:}
      \vspace{0.2em}
      \begin{itemize}
          \setlength{\itemsep}{0.25em}
        \item \textbf{Fortified \textit{Kastelli}:} Constant Aegean pirate raids (13th--18th c.) forced villages inland into defensive citadels (Pyrgos, Emporio). The windowless outer house walls formed an unbroken curtain wall.
        \item \textbf{Labyrinthine Wind Baffles:} Deliberately winding, asymmetrical street grids (e.g., Mykonos Chora) disoriented corsairs and physically broke gale-force \textit{Meltemi} wind tunnels.
        \item \textbf{Cistern Engineering (\textit{Sterna}):} Total lack of rivers required 100\% rainwater harvesting: every flat terrace and roof sloped into subterranean rock-hewn cisterns.
      \end{itemize}
    \end{column}
  \end{columns}
\end{frame}
